% =========================================================
% Academic Paper Template (Abstract-only or Full Paper)
% pdfLaTeX / Overleaf compatible
% =========================================================
\documentclass[11pt]{article}

% -------------------------
% Toggle: abstract-only vs full paper
% Set \abstractonlytrue for abstract submissions
% Set \abstractonlyfalse for full paper submissions
% -------------------------
\newif\ifabstractonly
\abstractonlyfalse
%\abstractonlytrue

% -------------------------
% Page layout & typography
% -------------------------
\usepackage[a4paper,margin=1in]{geometry}
\usepackage[T1]{fontenc}
\usepackage[english]{babel}
\usepackage{microtype}

% -------------------------
% Math, figures, tables
% -------------------------
\usepackage{amsmath,amssymb}
\usepackage{graphicx}
\usepackage{booktabs}
\usepackage{tabularx}
\usepackage{caption}
\usepackage{subcaption}

% -------------------------
% Lists and spacing
% -------------------------
\usepackage{enumitem}
\setlist[itemize]{leftmargin=*, itemsep=2pt, topsep=2pt}
\setlist[enumerate]{leftmargin=*, itemsep=2pt, topsep=2pt}

% -------------------------
% Links (conservative style)
% -------------------------
\usepackage[hidelinks]{hyperref}

% -------------------------
% Title block with affiliations
% -------------------------
\usepackage{authblk}

% -------------------------
% Optional: nicer section headings (conservative)
% -------------------------
\usepackage{titlesec}
\titleformat{\section}{\large\bfseries}{\thesection}{0.8em}{}
\titleformat{\subsection}{\normalsize\bfseries}{\thesubsection}{0.8em}{}
\titlespacing*{\section}{0pt}{10pt}{6pt}
\titlespacing*{\subsection}{0pt}{8pt}{4pt}

% -------------------------
% Bibliography (choose ONE approach)
% -------------------------
% Option A (recommended simple): BibTeX + plain bibliography
% \usepackage[numbers,sort&compress]{natbib}
% \bibliographystyle{plainnat}

% Option B: BibLaTeX (comment Option A above if you use this)
\usepackage[backend=bibtex,style=numeric,sorting=none]{biblatex}
\addbibresource{references.bib}

% =========================================================
% Paper metadata (placeholders)
% =========================================================
\title{\textbf{[Paper Title Goes Here: Concise, Descriptive, and Specific]}}

% Authors and affiliations (placeholders)
\author[1]{First Author}
\author[2]{Second Author}
\author[1,3]{Third Author}

\affil[1]{Department, Institution, City, Country\\
\texttt{first.author@institution.edu}}

\affil[2]{Department, Institution, City, Country\\
\texttt{second.author@institution.edu}}

\affil[3]{Department, Institution, City, Country\\
\texttt{third.author@institution.edu}}

\affil[*]{\textit{Corresponding author:} email@domain.com}

\date{} % leave empty for conference submissions unless required

% =========================================================
\begin{document}
\maketitle

% -------------------------
% Abstract + Keywords
% -------------------------
\begin{abstract}
\noindent
[Write a self-contained abstract here. Typical structure: (i) context/motivation, (ii) gap/problem, (iii) approach/methodology,
(iv) main results (quantitative if possible), (v) implications. Keep it within the conference word limit.]
\end{abstract}

\noindent\textbf{Keywords:} [keyword 1; keyword 2; keyword 3; keyword 4; keyword 5]

% -------------------------
% Abstract-only mode: stop after abstract if desired
% -------------------------
\ifabstractonly
\vspace{8pt}
\noindent\textbf{Submission type:} Abstract-only.

% Optional: add a short “extended abstract” section if the conference expects 1–2 pages
\section{Extended Abstract (Optional)}
[If required, include a brief description of the problem, key idea, and preliminary results. Keep concise.]

\printbibliography
\end{document}
\fi

% =========================================================
% Full paper sections (typical structure)
% =========================================================

\section{Introduction}
[Explain the problem context, why it matters, and the specific gap in existing work. Clearly state your contributions and briefly
summarize the paper structure.]

\subsection{Contributions}
[Use a short bullet list.]
\begin{itemize}
  \item Contribution 1: [What is new?]
  \item Contribution 2: [What is new?]
  \item Contribution 3: [What is new?]
\end{itemize}

\section{Related Work}
[Summarize prior work, group it by themes, and clearly position your work. Avoid a paper-by-paper list; synthesize.]

\section{Methodology}
[Describe assumptions, inputs/outputs, and the full pipeline. Provide enough detail for reproducibility.]

\subsection{Problem Definition}
[Formally define the problem; include symbols and notation.]

\subsection{Proposed Approach}
[Describe method components. Include equations, algorithms, or diagrams as needed.]

\subsection{Implementation Details}
[Software, hardware, training settings, hyperparameters, data splits, etc.]

\section{Experimental Setup}
[Describe datasets, evaluation metrics, baselines, and experimental protocol.]

\subsection{Datasets}
[Describe each dataset and how it is used.]

\subsection{Evaluation Metrics}
[Define metrics; include formulas if needed.]

\subsection{Baselines}
[Describe competitors/baselines used for comparison.]

\section{Results}
[Present main results with clear quantitative tables/figures. Report uncertainty if relevant.]

\subsection{Quantitative Results}
\begin{table}[ht]
  \centering
  \caption{[Placeholder caption: Main quantitative results (mean $\pm$ std).]}
  \begin{tabular}{lccc}
    \toprule
    Method & Metric 1 $\uparrow$ & Metric 2 $\downarrow$ & Runtime (ms) \\
    \midrule
    Baseline A & [ ] & [ ] & [ ] \\
    Baseline B & [ ] & [ ] & [ ] \\
    Proposed  & [ ] & [ ] & [ ] \\
    \bottomrule
  \end{tabular}
  \label{tab:main_results}
\end{table}

\subsection{Qualitative Results}
\begin{figure}[ht]
  \centering
  % Replace with your figure file
  \fbox{\includegraphics[width=0.9\linewidth]{figures/placeholder_figure.pdf}}
  \caption{[Placeholder caption: Qualitative comparison showing typical success cases and failure modes.]}
  \label{fig:qualitative}
\end{figure}

\section{Discussion}
[Interpret results, explain why your method works (or fails), and discuss practical implications, limitations, and edge cases.]

\subsection{Ablation Study (Optional)}
[Show impact of removing/altering components.]

\subsection{Limitations}
[Be explicit and constructive: where it fails, why, and how future work could address it.]

\section{Conclusion}
[Summarize key contributions and findings. Keep it brief. Mention future work.]

\section*{Acknowledgments}
[Funding sources, collaborators, institutional support. Follow conference guidelines.]

% =========================================================
% References
% =========================================================
\printbibliography

\end{document}
